\subsection*{CU-22: Invitar a un amigo a jugar}
\addcontentsline{toc}{subsection}{CU-22: Invitar a un amigo a jugar}
\label{cu-22}

\begin{fichacu}{CU-22}{Invitar a un amigo a jugar}
  \crudFila{Responsable}{Ángel Gabriel Aguilar Hernández}
  \crudFila{Actor principal}{Jugador que invita}
  \crudFila{Actores interesados}{Jugador invitado o destinatario del correo}
  \crudFila{Actores de sistema}{Servidor de partidas, Servicio de correo, Base de datos}
  \crudFila{Prototipos}{2a, 10a, 13a}
  \crudFila{Prioridad}{Must (debe estar en la entrega). Categoría (a) Obligatorio: característica 5 del cliente (invitación de amigos por correo o código).}
  \crudFila{Objetivo}{Permitir que el Jugador rete directamente a un amigo, o a cualquier jugador desde su menú, a una partida privada, sin que ninguno de los dos tenga que compartir un código; y que el anfitrión de una sala invite por correo a quien todavía no está en el juego, enviándole el código para entrar.}
  \crudFila{Disparador}{El Jugador selecciona ``invitar a jugar'' junto a un amigo en la \texttt{GUI\_\allowbreak{}Friends} o ``retar'' en el menú de un jugador; o, como anfitrión, selecciona ``Invitar por correo'' en la \texttt{GUI\_\allowbreak{}WaitingRoom} de su sala.}
  \textbf{Precondiciones} & \cuCelda{\begin{cureglas}
      \item \textbf{\mbox{PRE-01}} El Jugador tiene una \textit{Sesion} abierta y su token es válido.
      \item \textbf{\mbox{PRE-02}} El Jugador no participa en ninguna partida en línea sin terminar y no está en la \textit{ColaEmparejamiento} del \mbox{CU-08}.
      \item \textbf{\mbox{PRE-03}} El Servidor de partidas está en ejecución y tiene conexión disponible con la Base de datos.
    \end{cureglas}} \\ \hline
  \textbf{Flujo normal} & \cuCelda{\begin{cuenum}[start=1]
      \item El Jugador selecciona ``invitar a jugar'' junto a un amigo. \textit{(Ver \mbox{FA-10})}
      \item El SJM despliega la hoja rápida de invitación con el \textit{Modo} y el reloj de la última partida privada preseleccionados. \textit{(Ver \mbox{FA-01})}
      \item El Jugador confirma los ajustes y selecciona ``Enviar invitación''.
      \item El SJM envía el mensaje \texttt{INVITE\_\allowbreak{}TO\_\allowbreak{}MATCH\_\allowbreak{}REQUEST} con el \texttt{id\_\allowbreak{}usuario} del invitado, los ajustes y el token. \textit{(Ver \mbox{EX-03}, \mbox{EX-04})}
      \item El Servidor de partidas verifica que el token corresponda a una \textit{Sesion} abierta y comprueba la precondición \mbox{PRE-02}. \textit{(Ver \mbox{FA-02})}
      \item El Servidor de partidas verifica que el invitado exista y que no tenga la cuenta \texttt{ELIMINADA}. \textit{(Ver \mbox{FA-03})}
      \item El Servidor de partidas verifica que el invitado tenga una \textit{Sesion} abierta, que no participe en una partida en línea sin terminar y que no esté en la \textit{ColaEmparejamiento} del \mbox{CU-08} (\mbox{RN-05}). \textit{(Ver \mbox{FA-04}, \mbox{FA-05})}
      \item El Servidor de partidas verifica que no exista ya una \textit{Invitacion} pendiente del Jugador al mismo invitado (\mbox{RN-03}). \textit{(Ver \mbox{FA-06})}
    \end{cuenum}} \\
   & \cuCelda{\begin{cuenum}[start=9]
      \item Si el Jugador no está en una sala abierta, el Servidor de partidas crea la \textit{Sala} con los ajustes, conforme a los pasos 8 a 10 del \mbox{CU-09}; si ya está en una sala propia, usa esa. \textit{(Ver \mbox{EX-01})}
      \item El Servidor de partidas crea la \textit{Invitacion} de \texttt{canal} \texttt{JUEGO} con la \texttt{id\_\allowbreak{}sala}, el \texttt{id\_\allowbreak{}emisor}, el \texttt{id\_\allowbreak{}destinatario}, la \texttt{fecha\_\allowbreak{}invitacion}, la \texttt{fecha\_\allowbreak{}expiracion} a sesenta segundos y el \texttt{estado} igual a \texttt{PENDIENTE}. \textit{(Ver \mbox{EX-01})}
      \item El Servidor de partidas notifica al invitado con \texttt{MATCH\_\allowbreak{}INVITATION}, con el \texttt{nickname} del Jugador y los ajustes.
      \item El SJM del Jugador despliega la \texttt{GUI\_\allowbreak{}WaitingRoom} de la sala, con la plaza del invitado marcada como ``invitado, esperando respuesta''.
      \item El SJM del invitado muestra un aviso emergente con los ajustes y las opciones ``Aceptar'' y ``Rechazar'', con la cuenta atrás de los sesenta segundos. \textit{(Ver \mbox{FA-07}, \mbox{FA-08})}
      \item El invitado selecciona ``Aceptar'' y su SJM envía \texttt{ACCEPT\_\allowbreak{}INVITATION\_\allowbreak{}REQUEST}.
      \item El Servidor de partidas verifica que la \textit{Invitacion} siga \texttt{PENDIENTE}, no haya expirado y que la sala siga \texttt{ABIERTA} con plaza libre, y que el invitado no haya entrado mientras tanto en una partida o en la \textit{ColaEmparejamiento} del \mbox{CU-08} (\mbox{RN-05}). \textit{(Ver \mbox{FA-09}, \mbox{FA-14})}
      \item El Servidor de partidas marca la \textit{Invitacion} como \texttt{ACEPTADA} y agrega al invitado a la sala conforme a los pasos 9 a 11 del \mbox{CU-10}, sin pedir el código (\mbox{RN-04}). \textit{(Ver \mbox{EX-02})}
    \end{cuenum}} \\
   & \cuCelda{\begin{cuenum}[start=17]
      \item El SJM del invitado despliega la \texttt{GUI\_\allowbreak{}WaitingRoom}, y el flujo continúa en el paso 13 del \mbox{CU-10}.
      \item Fin del caso de uso.
    \end{cuenum}} \\ \hline
  \textbf{Flujos alternos} & \cuCelda{\textbf{\mbox{FA-01}: El Jugador reta a alguien que no es su amigo}\par
    \begin{cuenum}[start=1]
      \item El Jugador selecciona ``retar'' en el menú de un jugador desde el chat, el ranking o el fin de partida.
      \item El flujo continúa en el paso 2 del FN con las mismas comprobaciones; no hace falta amistad (\mbox{RN-02}).
    \end{cuenum}} \\
   & \cuCelda{\textbf{\mbox{FA-02}: El Jugador ya está en una partida o en la cola}\par
    \begin{cuenum}[start=1]
      \item El Servidor de partidas responde con \texttt{INVITATION\_\allowbreak{}FAILED} y el motivo \texttt{YA\_\allowbreak{}EN\_\allowbreak{}PARTIDA} o \texttt{YA\_\allowbreak{}EN\_\allowbreak{}COLA}.
      \item El SJM ofrece volver a la partida, que continúa en el \mbox{CU-18}, o cancelar la búsqueda del \mbox{CU-08}.
      \item Fin del caso de uso.
    \end{cuenum}} \\
   & \cuCelda{\textbf{\mbox{FA-03}: La cuenta del invitado no está disponible}\par
    \begin{cuenum}[start=1]
      \item El Servidor de partidas responde con \texttt{INVITATION\_\allowbreak{}FAILED} y el motivo \texttt{JUGADOR\_\allowbreak{}NO\_\allowbreak{}DISPONIBLE}.
      \item El SJM muestra que no es posible invitar a ese jugador.
      \item Fin del caso de uso.
    \end{cuenum}} \\
   & \cuCelda{\textbf{\mbox{FA-04}: El invitado no está conectado}\par
    \begin{cuenum}[start=1]
      \item El Servidor de partidas responde con \texttt{INVITATION\_\allowbreak{}FAILED} y el motivo \texttt{JUGADOR\_\allowbreak{}DESCONECTADO}, sin crear la sala ni la invitación.
      \item El SJM lo indica junto al amigo.
      \item Fin del caso de uso.
    \end{cuenum}} \\
   & \cuCelda{\textbf{\mbox{FA-05}: El invitado está jugando una partida o en la cola}\par
    \begin{cuenum}[start=1]
      \item El Servidor de partidas responde con \texttt{INVITATION\_\allowbreak{}FAILED} y el motivo \texttt{JUGADOR\_\allowbreak{}EN\_\allowbreak{}PARTIDA} o \texttt{JUGADOR\_\allowbreak{}EN\_\allowbreak{}COLA}, sin crear la sala ni la invitación.
      \item El SJM lo indica junto al amigo.
      \item Fin del caso de uso.
    \end{cuenum}} \\
   & \cuCelda{\textbf{\mbox{FA-06}: Ya hay una invitación pendiente al mismo jugador}\par
    \begin{cuenum}[start=1]
      \item El Servidor de partidas responde con \texttt{INVITATION\_\allowbreak{}ALREADY\_\allowbreak{}PENDING} y los segundos que le quedan.
      \item El SJM lleva al Jugador a la sala donde espera esa invitación.
      \item Fin del caso de uso.
    \end{cuenum}} \\
   & \cuCelda{\textbf{\mbox{FA-07}: El invitado rechaza}\par
    \begin{cuenum}[start=1]
      \item El invitado selecciona ``Rechazar'' y su SJM envía \texttt{DECLINE\_\allowbreak{}INVITATION\_\allowbreak{}REQUEST}.
      \item El Servidor de partidas marca la \textit{Invitacion} como \texttt{RECHAZADA} y notifica al Jugador con \texttt{INVITATION\_\allowbreak{}DECLINED}.
      \item El SJM del Jugador libera la plaza en la sala y lo indica; la sala sigue abierta para invitar a otro o compartir el código.
      \item Fin del caso de uso.
    \end{cuenum}} \\
   & \cuCelda{\textbf{\mbox{FA-08}: La invitación expira sin respuesta}\par
    \begin{cuenum}[start=1]
      \item Transcurren los sesenta segundos sin respuesta.
      \item El Servidor de partidas marca la \textit{Invitacion} como \texttt{EXPIRADA} y notifica a ambos.
      \item El SJM del invitado retira el aviso; el del Jugador libera la plaza.
      \item Fin del caso de uso.
    \end{cuenum}} \\
   & \cuCelda{\textbf{\mbox{FA-09}: La sala dejó de estar disponible antes de aceptar}\par
    \begin{cuenum}[start=1]
      \item El Servidor de partidas encuentra la sala cerrada, llena o en partida.
      \item El Servidor de partidas marca la \textit{Invitacion} como \texttt{EXPIRADA} y responde con \texttt{INVITATION\_\allowbreak{}NO\_\allowbreak{}LONGER\_\allowbreak{}VALID}.
      \item El SJM del invitado muestra que la invitación ya no es válida.
      \item Fin del caso de uso.
    \end{cuenum}} \\
   & \cuCelda{\textbf{\mbox{FA-10}: El anfitrión invita por correo desde su sala}\par
    \begin{cuenum}[start=1]
      \item Con su sala abierta en la \texttt{GUI\_\allowbreak{}WaitingRoom} (\mbox{CU-09}), el anfitrión selecciona ``Invitar por correo''. El SJM solo ofrece esta opción a un anfitrión con cuenta registrada (\mbox{RN-07}).
      \item El SJM despliega un campo para el correo del destinatario; el anfitrión lo captura y da click en ``Enviar''.
      \item El SJM valida que el texto tenga formato de dirección de correo. \textit{(Ver \mbox{FA-11})}
      \item El SJM envía el mensaje \texttt{INVITE\_\allowbreak{}BY\_\allowbreak{}EMAIL\_\allowbreak{}REQUEST} con el correo, el \texttt{id\_\allowbreak{}sala} y el token de sesión. \textit{(Ver \mbox{EX-03}, \mbox{EX-04})}
      \item El Servidor de partidas verifica que el token corresponda a una \textit{Sesion} abierta de una cuenta con \texttt{tipo\_\allowbreak{}cuenta} en \texttt{REGISTRADA}, que el Jugador sea el anfitrión de la \textit{Sala} y que esta siga \texttt{ABIERTA} (\mbox{RN-07}). \textit{(Ver \mbox{FA-13})}
      \item El Servidor de partidas vuelve a validar el formato del correo y cuenta las \textit{Invitacion} de \texttt{canal} \texttt{CORREO} que el Jugador creó en la última hora (\mbox{RN-06}). \textit{(Ver \mbox{FA-11}, \mbox{FA-12})}
      \item El Servidor de partidas solicita al Servicio de correo el envío al correo capturado de un mensaje en el \texttt{idioma\_\allowbreak{}preferido} del anfitrión, con su \texttt{nickname}, el \textit{Modo} de la sala, el \texttt{codigo} y las instrucciones para entrar con él, tenga o no cuenta el destinatario (\mbox{RN-08}, \mbox{RN-09}). \textit{(Ver \mbox{EX-05})}
      \item El Servidor de partidas crea la \textit{Invitacion} de \texttt{canal} \texttt{CORREO} con la \texttt{id\_\allowbreak{}sala}, el \texttt{id\_\allowbreak{}emisor}, el \texttt{correo\_\allowbreak{}destino} normalizado a minúsculas, la \texttt{fecha\_\allowbreak{}invitacion} y el \texttt{estado} igual a \texttt{ENVIADA}, sin \texttt{id\_\allowbreak{}destinatario} ni \texttt{fecha\_\allowbreak{}expiracion}. \textit{(Ver \mbox{EX-01})}
      \item El Servidor de partidas responde con \texttt{INVITE\_\allowbreak{}BY\_\allowbreak{}EMAIL\_\allowbreak{}OK}.
      \item El SJM indica junto al campo que la invitación se envió y lo vacía. La sala sigue abierta, y el destinatario entra cuando quiera con el código, conforme al \mbox{CU-10} (\mbox{RN-08}).
      \item Fin del caso de uso.
    \end{cuenum}} \\
   & \cuCelda{\textbf{\mbox{FA-11}: El correo capturado no tiene formato válido}\par
    \begin{cuenum}[start=1]
      \item El SJM escribe junto al campo que no es una dirección de correo válida, sin enviar nada al Servidor de partidas. Si quien lo rechaza es el Servidor de partidas, este registra la discrepancia en la bitácora del servidor y responde con \texttt{INVITE\_\allowbreak{}BY\_\allowbreak{}EMAIL\_\allowbreak{}FAILED} y el motivo \texttt{CORREO\_\allowbreak{}INVALIDO}.
      \item El flujo continúa en el paso 2 del \mbox{FA-10}.
    \end{cuenum}} \\
   & \cuCelda{\textbf{\mbox{FA-12}: Se alcanzó el límite de invitaciones por correo}\par
    \begin{cuenum}[start=1]
      \item El Servidor de partidas responde con \texttt{INVITE\_\allowbreak{}BY\_\allowbreak{}EMAIL\_\allowbreak{}FAILED}, el motivo \texttt{LIMITE\_\allowbreak{}ALCANZADO} y los minutos que faltan para poder invitar de nuevo (\mbox{RN-06}).
      \item El SJM lo indica junto al campo y deshabilita ``Invitar por correo'' durante ese tiempo; la sala sigue abierta y su código puede compartirse por otros medios.
      \item Fin del caso de uso.
    \end{cuenum}} \\
   & \cuCelda{\textbf{\mbox{FA-13}: La sala ya no está abierta o el Jugador no puede invitar por correo}\par
    \begin{cuenum}[start=1]
      \item El Servidor de partidas encuentra la sala cerrada o en partida, que el Jugador no es su anfitrión o que su cuenta es de invitado.
      \item El Servidor de partidas responde con \texttt{INVITE\_\allowbreak{}BY\_\allowbreak{}EMAIL\_\allowbreak{}FAILED} y el motivo \texttt{SALA\_\allowbreak{}NO\_\allowbreak{}DISPONIBLE} o \texttt{NO\_\allowbreak{}PERMITIDO}, sin solicitar ningún envío.
      \item El SJM muestra la \texttt{GUI\_\allowbreak{}MessageWarning} con el motivo y el estado actual de la sala.
      \item Fin del caso de uso.
    \end{cuenum}} \\
   & \cuCelda{\textbf{\mbox{FA-14}: El invitado entró en la cola o en una partida antes de aceptar}\par
    \begin{cuenum}[start=1]
      \item El Servidor de partidas responde al invitado con \texttt{INVITATION\_\allowbreak{}FAILED} y el motivo \texttt{YA\_\allowbreak{}EN\_\allowbreak{}COLA} o \texttt{YA\_\allowbreak{}EN\_\allowbreak{}PARTIDA}, sin marcar la \textit{Invitacion} (\mbox{RN-05}).
      \item El SJM del invitado ofrece cancelar la búsqueda del \mbox{CU-08} para poder aceptar mientras la invitación no expire; si la cancela, el flujo continúa en el paso 14 del FN.
      \item Si no la cancela, la invitación expira conforme al \mbox{FA-08}.
      \item Fin del caso de uso.
    \end{cuenum}} \\ \hline
  \textbf{Excepciones} & \cuCelda{\textbf{\mbox{EX-01}: Fallo al escribir en la base de datos}\par
    \begin{cuenum}[start=1]
      \item El Servidor de partidas revierte la transacción, de modo que no quede una \textit{Invitacion} hacia una sala que no se creó.
      \item El Servidor de partidas registra la excepción con un identificador de incidencia y responde con \texttt{SERVER\_\allowbreak{}ERROR}.
      \item El SJM muestra la \texttt{GUI\_\allowbreak{}MessageError} con el identificador. Si el fallo ocurrió al registrar una invitación por correo (paso 8 del \mbox{FA-10}), el mensaje ya se había entregado al Servicio de correo, y el SJM indica que el correo pudo haberse enviado, para que no se repita.
      \item Fin del caso de uso.
    \end{cuenum}} \\
   & \cuCelda{\textbf{\mbox{EX-02}: El invitado acepta y falla la entrada a la sala}\par
    \begin{cuenum}[start=1]
      \item El Servidor de partidas revierte la marca \texttt{ACEPTADA} junto con la entrada fallida.
      \item El SJM del invitado muestra el error y conserva el aviso mientras la invitación no expire.
      \item Fin del caso de uso.
    \end{cuenum}} \\
   & \cuCelda{\textbf{\mbox{EX-03}: Se agota el tiempo de espera de la respuesta}\par
    \begin{cuenum}[start=1]
      \item El SJM no reenvía la invitación y consulta el estado de la sala.
      \item El SJM muestra la sala con la invitación si se creó, o la lista de amigos si no. En el \mbox{FA-10}, el SJM permanece en la sala e indica que no pudo confirmar si el correo se envió.
    \end{cuenum}} \\
   & \cuCelda{\textbf{\mbox{EX-04}: La conexión se interrumpe}\par
    \begin{cuenum}[start=1]
      \item El SJM muestra la barra de conexión perdida.
      \item Si se pierde la conexión del invitado, la invitación expira conforme al \mbox{FA-08}; si se pierde la del Jugador, la sala se rige por el \mbox{EX-04} del \mbox{CU-09}.
    \end{cuenum}} \\
   & \cuCelda{\textbf{\mbox{EX-05}: El servicio de correo no está disponible}\par
    \begin{cuenum}[start=1]
      \item El Servidor de partidas agota los reintentos previstos hacia el Servicio de correo sin obtener confirmación de entrega.
      \item El Servidor de partidas no crea la \textit{Invitacion}, de modo que el intento fallido no cuenta para el límite del \mbox{RN-06}, y registra el fallo del servicio externo en la bitácora del servidor.
      \item El Servidor de partidas responde con \texttt{EXTERNAL\_\allowbreak{}SERVICE\_\allowbreak{}UNAVAILABLE}.
      \item El SJM muestra la \texttt{GUI\_\allowbreak{}MessageWarning} indicando que no fue posible enviar el correo y que, mientras tanto, puede compartir el código de la sala por su cuenta.
      \item El flujo continúa en el paso 2 del \mbox{FA-10}.
    \end{cuenum}} \\ \hline
  \textbf{Postcondiciones} & \cuCelda{\begin{cureglas}
      \item \textbf{\mbox{POST-01}} Si el caso termina por el FN, la \textit{Invitacion} está \texttt{ACEPTADA} y ambos jugadores son miembros de la misma \textit{Sala} abierta.
      \item \textbf{\mbox{POST-02}} Si el caso termina por el \mbox{FA-07}, el \mbox{FA-08} o el \mbox{FA-09}, la \textit{Invitacion} está \texttt{RECHAZADA} o \texttt{EXPIRADA} y el invitado no entró a la sala.
      \item \textbf{\mbox{POST-03}} No hay más de una \textit{Invitacion} pendiente del mismo emisor al mismo destinatario.
      \item \textbf{\mbox{POST-04}} La sala creada para la invitación se comporta como cualquier sala privada del \mbox{CU-09}, y la partida que salga de ella no modificará el elo ni la \textit{EstadisticaModo} de nadie.
      \item \textbf{\mbox{POST-05}} Si el caso termina por el \mbox{FA-10}, existe una \textit{Invitacion} de \texttt{canal} \texttt{CORREO} en estado \texttt{ENVIADA} con la sala, el emisor y el \texttt{correo\_\allowbreak{}destino}, y el Servicio de correo recibió el mensaje con el código. El destinatario no ocupa ninguna plaza mientras no entre con el \mbox{CU-10}.
    \end{cureglas}} \\ \hline
  \textbf{Reglas de negocio} & \cuCelda{\begin{cureglas}
      \item \textbf{\mbox{RN-01}} Una invitación a jugar crea o usa una sala privada del \mbox{CU-09}; no es un emparejamiento clasificatorio del \mbox{CU-08}. La partida que salga de ella es amistosa: no cambia el elo ni la \textit{EstadisticaModo} de nadie.
      \item \textbf{\mbox{RN-02}} Se puede retar a cualquier jugador, sea o no amigo.
      \item \textbf{\mbox{RN-03}} Una invitación dentro del juego expira a los sesenta segundos, y solo puede haber una pendiente del mismo emisor al mismo destinatario.
      \item \textbf{\mbox{RN-04}} Aceptar una invitación dentro del juego da acceso a la sala sin conocer su código.
      \item \textbf{\mbox{RN-05}} Solo se puede invitar dentro del juego a jugadores con sesión abierta que no estén en una partida en línea ni en la \textit{ColaEmparejamiento} del \mbox{CU-08}, y solo pueden aceptar mientras sigan así, porque un jugador no puede estar en una sala y en la cola a la vez.
    \end{cureglas}} \\
   & \cuCelda{\begin{cureglas}
      \item \textbf{\mbox{RN-06}} Un jugador puede enviar a lo sumo diez invitaciones por correo en una hora, para que el juego no sirva para enviar correo no deseado.
      \item \textbf{\mbox{RN-07}} Solo el anfitrión de una sala abierta, y solo con una cuenta registrada, puede invitar por correo: así cada envío queda atribuido a una cuenta con correo verificado.
      \item \textbf{\mbox{RN-08}} La invitación por correo solo lleva el código de la sala: no reserva plaza ni da acceso por sí misma. El destinatario entra con el \mbox{CU-10} como cualquier otro jugador, con una cuenta nueva del \mbox{CU-02} o como invitado del \mbox{CU-03}, y el código vale mientras la sala siga abierta.
      \item \textbf{\mbox{RN-09}} El correo capturado solo se usa para ese envío y para el límite del \mbox{RN-06}: no se asocia a ninguna cuenta ni se muestra a otros jugadores, y el mensaje no revela el correo del anfitrión. Va en el \texttt{idioma\_\allowbreak{}preferido} del anfitrión, porque el del destinatario no se conoce.
    \end{cureglas}} \\ \hline
  \textbf{Observación} & \cuCelda{\textit{Sobre por qué la invitación usa una sala en lugar de crear la partida directamente.}\par
    Crear la partida en cuanto el invitado acepta sería más rápido, pero dejaría dos formas de empezar una partida privada con reglas distintas. Pasar por la sala hace que la invitación sea solo una forma de entrar sin código: los ajustes, la marca de listo, el chat de espera y el inicio los resuelve el \mbox{CU-09}, y funciona igual para un reto uno contra uno que para completar una partida de cuatro invitando a varios amigos. Por la misma razón, la invitación por correo no necesita un mecanismo propio de entrada: envía el código y el destinatario entra con el \mbox{CU-10}.} \\ \hline
  \textbf{Observación} & \cuCelda{\textit{Sobre la cola y las plazas de sala en memoria (\mbox{D-24}).}\par
    La sala a la que se invita y sus plazas siguen la \mbox{D-24}: las plazas viven en memoria.} \\ \hline
  \textbf{Datos que toca} & \cuCelda{\begin{cudatos}
      \item \textbf{\textit{Sesion}:} lee la del emisor y comprueba la del destinatario \textit{(pasos 5, 7, \mbox{FA-10})}
      \item \textbf{\textit{Participacion}, \textit{ColaEmparejamiento}:} lee, para las exclusiones \textit{(pasos 5, 7, 15)}
      \item \textbf{\textit{Usuario}:} lee el invitado \textit{(paso 6)}; lee \texttt{tipo\_\allowbreak{}cuenta}, \texttt{nickname} e \texttt{idioma\_\allowbreak{}preferido} del anfitrión \textit{(\mbox{FA-10})}
      \item \textbf{\textit{Invitacion}:} lee las pendientes; crea; escribe \texttt{estado} \textit{(pasos 8, 10, 16, \mbox{FA-07}, \mbox{FA-08}, \mbox{FA-09})}
      \item \textbf{\textit{Invitacion}:} cuenta las de \texttt{canal} \texttt{CORREO} de la última hora; crea con \texttt{canal} \texttt{CORREO} \textit{(\mbox{FA-10})}
      \item \textbf{\textit{Sala}, \textit{SalaParticipante}:} crea o usa la sala; agrega al invitado \textit{(pasos 9, 16)}
      \item \textbf{\textit{Sala}, \textit{Modo}:} lee el anfitrión, el \texttt{estado}, el \texttt{codigo} y el modo de la sala \textit{(\mbox{FA-10})}
    \end{cudatos}} \\ \hline
\end{fichacu}
\clearpage
